\chapter{Неокейнсиански модел на малка отворена икономика}
\label{ch:model}

Настоящата глава извежда модела, който служи за структурна основа на изследването. Използва се вариантът на неокейнсианския модел за малка отворена икономика, представен в \textcite[гл.~7]{Gali2008} и разработен първоначално в \textcite{Gali_Monacelli2005}.

\section{Общи положения}

Световната икономика се моделира като континуум от малки отворени икономики, индексирани с $i \in [0,1]$. Всяка отделна икономика е с мярка нула, поради което поведението ѝ не влияе върху останалия свят. Икономиките се различават по реализациите на технологичните си шокове, но имат идентични предпочитания, технология и пазарна структура \parencite[с.~151]{Gali2008}.

За облекчаване на записа променливите без индекс се отнасят за разглежданата малка отворена икономика, променливите с индекс $i$ - за произволна друга икономика от континуума, а променливите със звездичка ($*$) - за световната икономика като цяло. Малките букви означават естествен логаритъм от съответните величини, т.е. $c_t := \ln C_t$, $p_t := \ln P_t$ и т.н.

\section{Домакинства}

\subsection{Предпочитания и агрегати на потреблението}

Представително домакинство в разглежданата икономика максимизира
\begin{equation}
\mathbb{E}_0 \sum_{t=0}^{\infty} \beta^t\, U(C_t, N_t),
\label{eq:hh_objective}
\end{equation}
където $N_t$ е количеството труд, а $C_t$ е съставен индекс (агрегатор) на потреблението с постоянна еластичност на заместване
\begin{equation}
C_t := \left[ (1-\alpha)^{\frac{1}{\eta}} \, C_{H,t}^{\frac{\eta-1}{\eta}}
+ \alpha^{\frac{1}{\eta}} \, C_{F,t}^{\frac{\eta-1}{\eta}} \right]^{\frac{\eta}{\eta-1}},
\qquad \alpha \in [0,1],\ \eta > 0.
\label{eq:C_index}
\end{equation}

Индексът на потребление на домашни стоки е
\begin{equation}
C_{H,t} := \left( \int_0^1 C_{H,t}(j)^{\frac{\varepsilon-1}{\varepsilon}}\, dj \right)^{\frac{\varepsilon}{\varepsilon-1}},
\qquad \varepsilon > 1,\ j \in [0,1],
\label{eq:CH_index}
\end{equation}
където $j$ индексира разновидностите, произвеждани в страната. Индексът на вносните стоки е
\begin{equation}
C_{F,t} := \left( \int_0^1 C_{i,t}^{\frac{\gamma-1}{\gamma}}\, di \right)^{\frac{\gamma}{\gamma-1}},
\qquad \gamma > 1,
\label{eq:CF_index}
\end{equation}
където $C_{i,t}$ е индекс на количеството стоки, внесени от страна $i$, зададен аналогично на \eqref{eq:CH_index}:
\begin{equation}
C_{i,t} := \left( \int_0^1 C_{i,t}(j)^{\frac{\varepsilon-1}{\varepsilon}}\, dj \right)^{\frac{\varepsilon}{\varepsilon-1}}.
\label{eq:Ci_index}
\end{equation}

Параметрите имат следната интерпретация \parencite[с.~152]{Gali2008}:
$\varepsilon$ е еластичността на заместване между разновидности, произведени в една и съща страна;
$\alpha$ измерва степента на отвореност на икономиката (съответно $1-\alpha$ е мярка за предпочитание към домашни стоки);
$\eta$ измерва заместимостта между домашни и чуждестранни стоки от гледна точка на домашния потребител;
$\gamma$ измерва заместимостта между стоки, произведени в различни чуждестранни страни.

Бюджетното ограничение има вида
\begin{equation}
\int_0^1 P_{H,t}(j) C_{H,t}(j)\, dj
+ \int_0^1 \!\! \int_0^1 P_{i,t}(j) C_{i,t}(j)\, dj\, di
+ \mathbb{E}_t\{ Q_{t,t+1} D_{t+1} \}
\le D_t + W_t N_t + T_t,
\label{eq:budget_full}
\end{equation}
където $P_{H,t}(j)$ е цената на домашната разновидност $j$, $P_{i,t}(j)$ е цената (в домашна валута) на разновидност $j$, внесена от страна $i$, $D_{t+1}$ е номиналното изплащане в период $t+1$ на портфейла, държан в края на период $t$, $W_t$ е номиналната работна заплата, $T_t$ са еднократни трансфери/данъци, а $Q_{t,t+1}$ е стохастичният дисконтов фактор за изплащания един период напред.

Приема се, че домакинствата имат достъп до пълен набор от условни вземания (contingent claims), търгувани на международно ниво. Освен това последователността от бюджетни ограничения се допълва с условие за платежоспособност от вида
\begin{equation}
\lim_{T \to \infty} \mathbb{E}_t\{ D_T \} \ge 0
\qquad \text{за всяко } t,
\label{eq:solvency}
\end{equation}
чиято интерпретация е, че се забраняват схеми на Понци и се осигурява изпълнението на условие за трансверсалност.

\subsection{Оптимално разпределение на разходите}

Задачата на домакинството се разлага на статична част (как да се разпредели даден размер разходи между стоките) и динамична част (колко да се потребява общо и колко да се работи). Статичната част се решава на три нива, като на всяко ниво задачата е от един и същ тип: минимизиране на разходите при фиксирана стойност на съответния агрегат. Показваме подробно първото ниво, след което същата процедура се повтаря.

\subsubsection*{Ниво 1: между разновидностите, произведени в една страна}

Домакинството решава
\begin{equation}
\min_{\{C_{H,t}(j)\}} \int_0^1 P_{H,t}(j) C_{H,t}(j)\, dj
\quad \text{при} \quad
\left( \int_0^1 C_{H,t}(j)^{\frac{\varepsilon-1}{\varepsilon}}\, dj \right)^{\frac{\varepsilon}{\varepsilon-1}} = C_{H,t}.
\label{eq:level1_problem}
\end{equation}
Полагаме $a := \frac{\varepsilon-1}{\varepsilon} \in (0,1)$, така че $\frac{1}{a} = \frac{\varepsilon}{\varepsilon-1}$ и ограничението се записва като $\left( \int_0^1 C_{H,t}(k)^a\, dk \right)^{1/a} = C_{H,t}$. Лагранжианът е
\begin{equation}
\mathcal{L} = \int_0^1 P_{H,t}(j) C_{H,t}(j)\, dj
+ \Lambda^{H}_t \left[ C_{H,t} - \left( \int_0^1 C_{H,t}(k)^{a}\, dk \right)^{\frac{1}{a}} \right].
\label{eq:level1_lagrangian}
\end{equation}

Производната на първия член по $C_{H,t}(j)$ е $P_{H,t}(j)$, а производната на $\Lambda^{H}_t C_{H,t}$ е нула, тъй като $C_{H,t}$ е зададено. За производната на агрегата имаме
\begin{equation}
\frac{\partial}{\partial C_{H,t}(j)} \left( \int_0^1 C_{H,t}(k)^{a}\, dk \right)^{\frac{1}{a}}
= \frac{1}{a} \left( \int_0^1 C_{H,t}(k)^{a}\, dk \right)^{\frac{1}{a}-1} \cdot a\, C_{H,t}(j)^{a-1}.
\label{eq:level1_deriv}
\end{equation}
Означавайки $X_t := \int_0^1 C_{H,t}(k)^{a}\, dk$, от ограничението следва $X_t^{1/a} = C_{H,t}$, т.е. $X_t = C_{H,t}^{a}$. Тогава
\[
X_t^{\frac{1}{a}-1} = \left( C_{H,t}^{a} \right)^{\frac{1}{a}-1} = C_{H,t}^{\,1-a} = C_{H,t}^{\,1/\varepsilon},
\qquad
a - 1 = -\frac{1}{\varepsilon},
\]
поради което \eqref{eq:level1_deriv} се свежда до $C_{H,t}^{1/\varepsilon} C_{H,t}(j)^{-1/\varepsilon}$. Необходимото условие за оптималност е
\begin{equation}
\frac{\partial \mathcal{L}}{\partial C_{H,t}(j)}
= P_{H,t}(j) - \Lambda^{H}_t\, C_{H,t}^{\frac{1}{\varepsilon}}\, C_{H,t}(j)^{-\frac{1}{\varepsilon}} = 0,
\label{eq:level1_foc}
\end{equation}
откъдето
\begin{equation}
C_{H,t}(j) = \left( \frac{P_{H,t}(j)}{\Lambda^{H}_t} \right)^{-\varepsilon} C_{H,t}.
\label{eq:level1_demand_lambda}
\end{equation}

За да се определи множителят $\Lambda^{H}_t$, умножаваме \eqref{eq:level1_demand_lambda} по $P_{H,t}(j)$ и интегрираме:
\begin{equation}
\int_0^1 P_{H,t}(j) C_{H,t}(j)\, dj
= \left( \Lambda^{H}_t \right)^{\varepsilon} C_{H,t} \int_0^1 P_{H,t}(j)^{1-\varepsilon}\, dj.
\label{eq:level1_expenditure}
\end{equation}
Дефинираме ценовите индекси
\begin{equation}
P_{H,t} := \left( \int_0^1 P_{H,t}(j)^{1-\varepsilon}\, dj \right)^{\frac{1}{1-\varepsilon}},
\qquad
P_{i,t} := \left( \int_0^1 P_{i,t}(j)^{1-\varepsilon}\, dj \right)^{\frac{1}{1-\varepsilon}},
\label{eq:PH_index}
\end{equation}
така че $\int_0^1 P_{H,t}(j)^{1-\varepsilon}\, dj = P_{H,t}^{1-\varepsilon}$. Заместването на \eqref{eq:level1_demand_lambda} в ограничението дава
\[
\left( \int_0^1 \left[ \left( \tfrac{P_{H,t}(k)}{\Lambda^{H}_t} \right)^{-\varepsilon} C_{H,t} \right]^{a} dk \right)^{\frac{1}{a}}
= C_{H,t} \left( \left( \Lambda^{H}_t \right)^{\varepsilon - 1} P_{H,t}^{1-\varepsilon} \right)^{\frac{1}{a}} \overset{!}{=} C_{H,t},
\]
което е изпълнено само при $\Lambda^{H}_t = P_{H,t}$. Следователно
\begin{equation}
C_{H,t}(j) = \left( \frac{P_{H,t}(j)}{P_{H,t}} \right)^{-\varepsilon} C_{H,t};
\qquad
C_{i,t}(j) = \left( \frac{P_{i,t}(j)}{P_{i,t}} \right)^{-\varepsilon} C_{i,t},
\label{eq:demand_variety}
\end{equation}
а от \eqref{eq:level1_expenditure} следва връзката
\begin{equation}
\int_0^1 P_{H,t}(j) C_{H,t}(j)\, dj = P_{H,t} C_{H,t},
\qquad
\int_0^1 P_{i,t}(j) C_{i,t}(j)\, dj = P_{i,t} C_{i,t}.
\label{eq:expenditure_variety}
\end{equation}

Съществено е, че множителят на Лагранж съвпада с ценовия индекс. Това е общо свойство на CES-агрегаторите и се повтаря на всяко от следващите нива.

\subsubsection*{Ниво 2: между страните на произход на вноса}

Прилагайки същата процедура към \eqref{eq:CF_index} с $b := \frac{\gamma-1}{\gamma}$ и множител $\Lambda^{F}_t$, получаваме необходимото условие
\[
P_{i,t} - \Lambda^{F}_t\, C_{F,t}^{\frac{1}{\gamma}} C_{i,t}^{-\frac{1}{\gamma}} = 0
\quad \Longrightarrow \quad
C_{i,t} = \left( \frac{P_{i,t}}{\Lambda^{F}_t} \right)^{-\gamma} C_{F,t},
\]
а определянето на множителя по същия начин дава $\Lambda^{F}_t = P_{F,t}$, където
\begin{equation}
P_{F,t} := \left( \int_0^1 P_{i,t}^{1-\gamma}\, di \right)^{\frac{1}{1-\gamma}}.
\label{eq:PF_index}
\end{equation}
Следователно
\begin{equation}
C_{i,t} = \left( \frac{P_{i,t}}{P_{F,t}} \right)^{-\gamma} C_{F,t},
\qquad
\int_0^1 P_{i,t} C_{i,t}\, di = P_{F,t} C_{F,t}.
\label{eq:demand_country}
\end{equation}

\subsubsection*{Ниво 3: между домашни и вносни стоки}

На третото ниво агрегаторът \eqref{eq:C_index} съдържа тегла, поради което извеждането се различава по константите. Задачата е
\[
\min_{C_{H,t},\, C_{F,t}} P_{H,t} C_{H,t} + P_{F,t} C_{F,t}
\quad \text{при} \quad
\left[ (1-\alpha)^{\frac{1}{\eta}} C_{H,t}^{\frac{\eta-1}{\eta}} + \alpha^{\frac{1}{\eta}} C_{F,t}^{\frac{\eta-1}{\eta}} \right]^{\frac{\eta}{\eta-1}} = C_t.
\]
Означавайки съдържанието на квадратните скоби с $\Xi_t$ и множителя на Лагранж с $\Lambda^{C}_t$, диференцирането по $C_{H,t}$ дава
\begin{equation}
P_{H,t} = \Lambda^{C}_t \cdot \frac{\eta}{\eta-1}\, \Xi_t^{\frac{\eta}{\eta-1}-1} \cdot (1-\alpha)^{\frac{1}{\eta}} \cdot \frac{\eta-1}{\eta}\, C_{H,t}^{-\frac{1}{\eta}}
= \Lambda^{C}_t\, \Xi_t^{\frac{1}{\eta-1}} (1-\alpha)^{\frac{1}{\eta}} C_{H,t}^{-\frac{1}{\eta}}.
\label{eq:level3_foc}
\end{equation}
Тъй като $\Xi_t^{\eta/(\eta-1)} = C_t$, е в сила $\Xi_t^{1/(\eta-1)} = C_t^{1/\eta}$, откъдето
\begin{equation}
C_{H,t} = (1-\alpha) \left( \frac{P_{H,t}}{\Lambda^{C}_t} \right)^{-\eta} C_t,
\qquad
C_{F,t} = \alpha \left( \frac{P_{F,t}}{\Lambda^{C}_t} \right)^{-\eta} C_t.
\label{eq:level3_demand_lambda}
\end{equation}

За определяне на $\Lambda^{C}_t$ заместваме \eqref{eq:level3_demand_lambda} в агрегатора. Първият член дава
\[
(1-\alpha)^{\frac{1}{\eta}} \left[ (1-\alpha) \left( \tfrac{P_{H,t}}{\Lambda^{C}_t} \right)^{-\eta} C_t \right]^{\frac{\eta-1}{\eta}}
= (1-\alpha) \left( \frac{P_{H,t}}{\Lambda^{C}_t} \right)^{1-\eta} C_t^{\frac{\eta-1}{\eta}},
\]
и аналогично за втория. Сумата е
\[
\Xi_t = C_t^{\frac{\eta-1}{\eta}} \left( \Lambda^{C}_t \right)^{\eta-1} \left[ (1-\alpha) P_{H,t}^{1-\eta} + \alpha P_{F,t}^{1-\eta} \right].
\]
Ограничението изисква $\Xi_t = C_t^{(\eta-1)/\eta}$, което дава $\left( \Lambda^{C}_t \right)^{\eta-1}\left[ (1-\alpha) P_{H,t}^{1-\eta} + \alpha P_{F,t}^{1-\eta} \right] = 1$. Дефинирайки индекса на потребителските цени (ИПЦ)
\begin{equation}
P_t := \left[ (1-\alpha) P_{H,t}^{1-\eta} + \alpha P_{F,t}^{1-\eta} \right]^{\frac{1}{1-\eta}},
\label{eq:CPI}
\end{equation}
условието се записва като $\left( \Lambda^{C}_t \right)^{\eta-1} P_t^{1-\eta} = 1$, откъдето $\Lambda^{C}_t = P_t$. Следователно
\begin{equation}
C_{H,t} = (1-\alpha) \left( \frac{P_{H,t}}{P_t} \right)^{-\eta} C_t;
\qquad
C_{F,t} = \alpha \left( \frac{P_{F,t}}{P_t} \right)^{-\eta} C_t.
\label{eq:demand_HF}
\end{equation}

При $\eta = 1$ или при равни ценови индекси за домашни и чуждестранни стоки параметърът $\alpha$ съответства точно на дела на разходите за внос в общото потребление, което оправдава интерпретацията му като индекс на отвореност \parencite[с.~153]{Gali2008}.

\subsubsection*{Свиване на бюджетното ограничение}

Общите разходи за потребление се получават от \eqref{eq:demand_HF}:
\begin{align}
P_{H,t} C_{H,t} + P_{F,t} C_{F,t}
&= (1-\alpha) P_{H,t}^{1-\eta} P_t^{\eta} C_t + \alpha P_{F,t}^{1-\eta} P_t^{\eta} C_t \notag \\
&= P_t^{\eta} C_t \left[ (1-\alpha) P_{H,t}^{1-\eta} + \alpha P_{F,t}^{1-\eta} \right]
= P_t^{\eta} C_t\, P_t^{1-\eta} = P_t C_t.
\label{eq:total_expenditure}
\end{align}

Комбинирайки \eqref{eq:total_expenditure} с \eqref{eq:expenditure_variety} и \eqref{eq:demand_country}, лявата страна на \eqref{eq:budget_full} се свежда до $P_t C_t + \mathbb{E}_t\{Q_{t,t+1}D_{t+1}\}$, т.е. бюджетното ограничение приема познатия вид
\begin{equation}
P_t C_t + \mathbb{E}_t\{ Q_{t,t+1} D_{t+1} \} \le D_t + W_t N_t + T_t.
\label{eq:budget_compact}
\end{equation}
Така след агрегиране на трите нива задачата на домакинството формално съвпада със задачата от затворената икономика, като единствената разлика е, че $P_t$ вече е ИПЦ, а не цената на единствената домашна стока.

\subsection{Условия за оптималност}

Периодната функция на полезност се специализира във вида
\begin{equation}
U(C_t, N_t) = \frac{C_t^{1-\sigma}}{1-\sigma} - \frac{N_t^{1+\varphi}}{1+\varphi},
\qquad \sigma > 0,\ \varphi \ge 0,
\label{eq:utility}
\end{equation}
за която $U_{c,t} = C_t^{-\sigma}$ и $U_{n,t} = -N_t^{\varphi}$.

\subsubsection*{Вътрепериодно условие}

Изравняването на пределната норма на заместване между потребление и труд с реалната работна заплата дава
\begin{equation}
-\frac{U_{n,t}}{U_{c,t}} = \frac{N_t^{\varphi}}{C_t^{-\sigma}} = C_t^{\sigma} N_t^{\varphi} = \frac{W_t}{P_t}.
\label{eq:labor_supply}
\end{equation}
Логаритмуването на \eqref{eq:labor_supply} е точно (не се изисква приближение) и дава
\begin{equation}
w_t - p_t = \sigma c_t + \varphi n_t.
\label{eq:labor_supply_log}
\end{equation}

\subsubsection*{Междувременно условие}

За извеждане на междувременното условие се използва фактът, че при пълни пазари за всяко състояние на природата трябва да е в сила
\begin{equation}
\frac{V_{t,t+1}}{P_t} C_t^{-\sigma} = \xi_{t,t+1}\, \beta\, C_{t+1}^{-\sigma} \frac{1}{P_{t+1}},
\label{eq:arrow}
\end{equation}
където $V_{t,t+1}$ е цената в период $t$ (в домашна валута) на ценна книга на Ароу (Arrow security), изплащаща единица домашна валута при реализиране на конкретно състояние в период $t+1$, а $\xi_{t,t+1}$ е вероятността за реализиране на това състояние. Лявата страна измерва загубата на полезност от намаленото текущо потребление, необходимо за закупуване на книгата, а дясната - очакваната печалба от допълнителното потребление, което изплащането прави възможно.

Тъй като $Q_{t,t+1} := V_{t,t+1} / \xi_{t,t+1}$, разделянето на \eqref{eq:arrow} на $\xi_{t,t+1}$ и на $C_t^{-\sigma}/P_t$ дава
\begin{equation}
Q_{t,t+1} = \beta \left( \frac{C_{t+1}}{C_t} \right)^{-\sigma} \left( \frac{P_t}{P_{t+1}} \right),
\label{eq:sdf}
\end{equation}
което трябва да е изпълнено за всяко състояние в $t$ и $t+1$. Вземайки условно очакване от двете страни и полагайки $Q_t := \mathbb{E}_t\{Q_{t,t+1}\}$, получаваме уравнението на Ойлер
\begin{equation}
Q_t = \beta\, \mathbb{E}_t \left\{ \left( \frac{C_{t+1}}{C_t} \right)^{-\sigma} \left( \frac{P_t}{P_{t+1}} \right) \right\}.
\label{eq:euler}
\end{equation}

\subsubsection*{Линеаризация на уравнението на Ойлер}

Въвеждаме означенията $i_t := -\ln Q_t$, $\rho := -\ln \beta$ и $\pi_{t+1} := p_{t+1} - p_t$. Тогава
\[
Q_t = e^{-i_t}, \qquad \beta = e^{-\rho}, \qquad
\left( \frac{C_{t+1}}{C_t} \right)^{-\sigma} = e^{-\sigma \Delta c_{t+1}}, \qquad
\frac{P_t}{P_{t+1}} = e^{-\pi_{t+1}},
\]
и \eqref{eq:euler} приема вида $e^{-i_t} = \mathbb{E}_t\{ e^{-\rho - \sigma\Delta c_{t+1} - \pi_{t+1}} \}$. Умножавайки двете страни по $e^{\,i_t}$, получаваме
\begin{equation}
1 = \mathbb{E}_t \left\{ e^{\, X_{t+1}} \right\},
\qquad X_{t+1} := i_t - \sigma \Delta c_{t+1} - \pi_{t+1} - \rho.
\label{eq:euler_exp}
\end{equation}

Разглеждаме равновесие, при което потреблението расте с постоянен темп $\Delta c_t = g$, а инфлацията е постоянна, $\pi_t = \pi$. Тогава $X = i - \sigma g - \pi - \rho$ и от \eqref{eq:euler_exp} следва $e^{EQ} := e^{X} = 1$, т.е. $X = 0$. Линеаризирайки $e^{X_{t+1}}$ около това равновесие, имаме
\begin{equation}
e^{X_{t+1}} \approx e^{EQ} + e^{EQ}(X_{t+1} - X) = 1 + X_{t+1}.
\label{eq:euler_linearize}
\end{equation}
Заместването в \eqref{eq:euler_exp} дава $1 = \mathbb{E}_t\{1 + X_{t+1}\}$, т.е. $\mathbb{E}_t\{X_{t+1}\} = 0$:
\[
i_t - \sigma\, \mathbb{E}_t\{ \Delta c_{t+1} \} - \mathbb{E}_t\{\pi_{t+1}\} - \rho = 0.
\]
Решавайки спрямо $\mathbb{E}_t\{\Delta c_{t+1}\} = \mathbb{E}_t\{c_{t+1}\} - c_t$, получаваме лог-линеаризираното уравнение на Ойлер
\begin{equation}
c_t = \mathbb{E}_t\{ c_{t+1} \} - \frac{1}{\sigma} \left( i_t - \mathbb{E}_t\{ \pi_{t+1} \} - \rho \right).
\label{eq:euler_log}
\end{equation}

Тук $i_t$ е краткосрочният номинален лихвен процент, $\rho$ е нормата на времево предпочитание, а $\pi_t$ е \emph{ИПЦ инфлацията} \parencite[с.~155]{Gali2008}. Уравнението \eqref{eq:euler_log} е записано спрямо ИПЦ инфлацията, а не спрямо домашната инфлация; преминаването между двете изисква въвеждането на условията на търговия.

\section{Относителни цени и международни връзки}

\subsection{Условия на търговия и ИПЦ инфлация}

Двустранните условия на търговия (terms of trade) със страна $i$ се дефинират като $S_{i,t} := P_{i,t} / P_{H,t}$, т.е. цената на стоките на страна $i$, изразена в единици домашни стоки. Ефективните условия на търговия са
\begin{equation}
S_t := \frac{P_{F,t}}{P_{H,t}} = \left( \int_0^1 S_{i,t}^{1-\gamma}\, di \right)^{\frac{1}{1-\gamma}},
\label{eq:tot_def}
\end{equation}
където второто равенство следва от \eqref{eq:PF_index} след разделяне на $P_{H,t}^{1-\gamma}$.

\paragraph{Лог-линеаризация на \eqref{eq:tot_def}.}
Логаритмувайки и записвайки $S_{i,t} = e^{s_{i,t}}$, имаме
\[
s_t = \frac{1}{1-\gamma} \ln \int_0^1 e^{(1-\gamma) s_{i,t}}\, di.
\]
Около симетрично стационарно състояние с $S_{i,t} = 1$ (т.е. $s_{i,t} = 0$) за всяко $i$ приближението от първи ред на подинтегралната функция е $e^{(1-\gamma)s_{i,t}} \approx 1 + (1-\gamma)s_{i,t}$, откъдето
\[
\int_0^1 e^{(1-\gamma) s_{i,t}}\, di \approx 1 + (1-\gamma) \int_0^1 s_{i,t}\, di.
\]
Използвайки $\ln(1+x) \approx x$, получаваме
\begin{equation}
s_t \approx \frac{1}{1-\gamma} \cdot (1-\gamma) \int_0^1 s_{i,t}\, di = \int_0^1 s_{i,t}\, di,
\qquad s_t := \ln S_t = p_{F,t} - p_{H,t}.
\label{eq:tot_loglin}
\end{equation}
Връзката е точна при $\gamma = 1$, тъй като тогава \eqref{eq:tot_def} се свежда до геометрична средна.

\paragraph{Лог-линеаризация на ИПЦ.}
Разделяйки \eqref{eq:CPI} на $P_{H,t}$, получаваме
\begin{equation}
\frac{P_t}{P_{H,t}} = \left[ (1-\alpha) + \alpha S_t^{1-\eta} \right]^{\frac{1}{1-\eta}},
\label{eq:cpi_ratio}
\end{equation}
т.е. $p_t - p_{H,t} = \frac{1}{1-\eta} \ln\!\left[ (1-\alpha) + \alpha e^{(1-\eta)s_t} \right] =: f(s_t)$. При $s_t = 0$ имаме $f(0) = \frac{1}{1-\eta}\ln 1 = 0$, а производната е
\[
f'(s_t) = \frac{1}{1-\eta} \cdot \frac{\alpha(1-\eta)e^{(1-\eta)s_t}}{(1-\alpha) + \alpha e^{(1-\eta)s_t}},
\qquad
f'(0) = \frac{\alpha}{(1-\alpha)+\alpha} = \alpha.
\]
Следователно приближението от първи ред е
\begin{equation}
p_t = (1-\alpha) p_{H,t} + \alpha p_{F,t} = p_{H,t} + \alpha s_t,
\label{eq:cpi_loglin}
\end{equation}
което е точно при $\eta = 1$, тъй като тогава $P_t = P_{H,t}^{1-\alpha} P_{F,t}^{\alpha}$ \parencite[с.~155]{Gali2008}.

Дефинирайки домашната инфлация като темпа на изменение на индекса на цените на домашно произведените стоки, $\pi_{H,t} := p_{H,t} - p_{H,t-1}$, и вземайки първа разлика на \eqref{eq:cpi_loglin}, получаваме
\begin{equation}
\pi_t = \pi_{H,t} + \alpha\, \Delta s_t.
\label{eq:pi_cpi_domestic}
\end{equation}
Разликата между двете мерки за инфлация е пропорционална на изменението на условията на търговия, с коефициент на пропорционалност степента на отвореност $\alpha$. При $\alpha = 0$ двете мерки съвпадат.

\subsection{Закон за единната цена и реален обменен курс}

Приема се, че законът за единната цена (law of one price) е в сила за всяка отделна стока и във всеки момент:
\begin{equation}
P_{i,t}(j) = \mathcal{E}_{i,t}\, P^i_{i,t}(j), \qquad \forall\, i, j \in [0,1],
\label{eq:lop}
\end{equation}
където $\mathcal{E}_{i,t}$ е двустранният номинален обменен курс (цената на валутата на страна $i$ в единици домашна валута), а $P^i_{i,t}(j)$ е цената на стоката, изразена във валутата на страна $i$.

Заместването на \eqref{eq:lop} в дефиницията на $P_{i,t}$ от \eqref{eq:PH_index} дава
\[
P_{i,t} = \left( \int_0^1 \left[ \mathcal{E}_{i,t} P^i_{i,t}(j) \right]^{1-\varepsilon} dj \right)^{\frac{1}{1-\varepsilon}}
= \mathcal{E}_{i,t} \left( \int_0^1 P^i_{i,t}(j)^{1-\varepsilon} dj \right)^{\frac{1}{1-\varepsilon}}
= \mathcal{E}_{i,t} P^i_{i,t},
\]
тъй като $\mathcal{E}_{i,t}$ не зависи от $j$ и може да бъде изнесено пред интеграла. Лог-линеаризирайки дефиницията \eqref{eq:PF_index} по същия начин както \eqref{eq:tot_def}, получаваме
\begin{equation}
p_{F,t} = \int_0^1 \left( e_{i,t} + p^i_{i,t} \right) di = e_t + p^*_t,
\label{eq:pF_loglin}
\end{equation}
където $e_t := \int_0^1 e_{i,t}\, di$ е (логаритъмът на) ефективният номинален обменен курс, а $p^*_t := \int_0^1 p^i_{i,t}\, di$ е световният ценови индекс. За света като цяло няма разлика между ИПЦ и индекса на домашните цени.

Комбинирайки \eqref{eq:pF_loglin} с $s_t = p_{F,t} - p_{H,t}$, получаваме
\begin{equation}
s_t = e_t + p^*_t - p_{H,t}.
\label{eq:tot_exchange}
\end{equation}

Двустранният реален обменен курс се дефинира като $\mathcal{Q}_{i,t} := \mathcal{E}_{i,t} P^i_t / P_t$, т.е. отношението на ИПЦ на двете страни, изразени в домашна валута. За ефективния реален обменен курс $q_t := \int_0^1 q_{i,t}\, di$ последователно имаме
\begin{align}
q_t &= \int_0^1 \left( e_{i,t} + p^i_t - p_t \right) di
= e_t + p^*_t - p_t \notag \\
&= \underbrace{(s_t + p_{H,t})}_{\text{от }\eqref{eq:tot_exchange}} - p_t
= s_t - \underbrace{(p_t - p_{H,t})}_{= \alpha s_t \ \text{от}\ \eqref{eq:cpi_loglin}}
= (1-\alpha) s_t,
\label{eq:rer}
\end{align}
където последното равенство е точно при $\eta = 1$, а в общия случай - с точност до първи ред \parencite[с.~156]{Gali2008}.

\subsection{Международно споделяне на риска}

При пълни международни финансови пазари условие, аналогично на \eqref{eq:arrow}, е в сила и за представителното домакинство във всяка друга страна $i$. Записано чрез стохастичния дисконтов фактор на разглежданата икономика, то има вида
\begin{equation}
\beta \left( \frac{C^i_{t+1}}{C^i_t} \right)^{-\sigma}
\left( \frac{P^i_t}{P^i_{t+1}} \right)
\left( \frac{\mathcal{E}_{i,t}}{\mathcal{E}_{i,t+1}} \right) = Q_{t,t+1}.
\label{eq:arrow_foreign}
\end{equation}
Наличието на членовете с обменния курс отразява факта, че ценната книга се купува на цена и с изплащане, деноминирани във валутата на разглежданата икономика, и следователно трябва да бъде преизчислена във валутата на страна $i$.

\paragraph{Извеждане на условието за споделяне на риска.}
Приравнявайки десните страни на \eqref{eq:sdf} и \eqref{eq:arrow_foreign} и съкращавайки $\beta$, получаваме
\[
\left( \frac{C_{t+1}}{C_t} \right)^{-\sigma} \frac{P_t}{P_{t+1}}
= \left( \frac{C^i_{t+1}}{C^i_t} \right)^{-\sigma} \frac{P^i_t}{P^i_{t+1}} \cdot \frac{\mathcal{E}_{i,t}}{\mathcal{E}_{i,t+1}}.
\]
Прегрупирайки така, че величините за домашната икономика да бъдат отляво, имаме
\begin{equation}
\left( \frac{C_{t+1}/C^i_{t+1}}{C_t/C^i_t} \right)^{-\sigma}
= \frac{P_{t+1}}{P_t} \cdot \frac{P^i_t}{P^i_{t+1}} \cdot \frac{\mathcal{E}_{i,t}}{\mathcal{E}_{i,t+1}}.
\label{eq:rs_step1}
\end{equation}
От дефиницията $\mathcal{Q}_{i,t} = \mathcal{E}_{i,t}P^i_t/P_t$ следва
\[
\frac{\mathcal{Q}_{i,t+1}}{\mathcal{Q}_{i,t}}
= \frac{\mathcal{E}_{i,t+1}}{\mathcal{E}_{i,t}} \cdot \frac{P^i_{t+1}}{P^i_t} \cdot \frac{P_t}{P_{t+1}},
\]
т.е. дясната страна на \eqref{eq:rs_step1} е точно $\mathcal{Q}_{i,t}/\mathcal{Q}_{i,t+1}$. Следователно
\begin{equation}
\left( \frac{C_{t+1}}{C^i_{t+1}} \right)^{-\sigma} \mathcal{Q}_{i,t+1}
= \left( \frac{C_t}{C^i_t} \right)^{-\sigma} \mathcal{Q}_{i,t},
\label{eq:rs_step2}
\end{equation}
т.е. величината $(C_t/C^i_t)^{-\sigma}\mathcal{Q}_{i,t}$ е постоянна във времето. Оттук
\begin{equation}
C_t = \vartheta_i\, C^i_t\, \mathcal{Q}_{i,t}^{\frac{1}{\sigma}},
\label{eq:risk_sharing_level}
\end{equation}
където $\vartheta_i$ е константа, зависеща от началните относителни нетни активни позиции. Приемайки симетрични начални условия (нулеви нетни чуждестранни активи), имаме $\vartheta_i = 1$ за всяко $i$.

Логаритмувайки \eqref{eq:risk_sharing_level} при $\vartheta_i = 1$, получаваме $c_t = c^i_t + \frac{1}{\sigma} q_{i,t}$. Интегрирайки по $i$ и използвайки, че лявата страна не зависи от $i$:
\begin{equation}
c_t = \int_0^1 c^i_t\, di + \frac{1}{\sigma} \int_0^1 q_{i,t}\, di
= c^*_t + \frac{1}{\sigma} q_t
= c^*_t + \left( \frac{1-\alpha}{\sigma} \right) s_t,
\label{eq:risk_sharing}
\end{equation}
където $c^*_t := \int_0^1 c^i_t\, di$ е индексът на световното потребление, а последното равенство използва \eqref{eq:rer} \parencite[с.~157]{Gali2008}. Условието \eqref{eq:risk_sharing} заменя необходимостта от проследяване на динамиката на нетните чуждестранни активи с статична връзка между домашното потребление, световното потребление и относителните цени.

\subsection{Непокрит лихвен паритет}

При пълни международни финансови пазари равновесната цена (в домашна валута) на безрискова облигация, деноминирана във валутата на страна $i$, е
\[
\mathcal{E}_{i,t} Q^i_t = \mathbb{E}_t\{ Q_{t,t+1} \mathcal{E}_{i,t+1} \}
\quad \Longleftrightarrow \quad
Q^i_t = \mathbb{E}_t\left\{ Q_{t,t+1} \frac{\mathcal{E}_{i,t+1}}{\mathcal{E}_{i,t}} \right\}.
\]
Записвайки $Q^i_t = e^{-i^i_t}$ и умножавайки по $e^{\,i^i_t}$, получаваме $\mathbb{E}_t\{Q_{t,t+1} e^{\,i^i_t}(\mathcal{E}_{i,t+1}/\mathcal{E}_{i,t})\} = 1$. Аналогично от домашното условие $Q_t = \mathbb{E}_t\{Q_{t,t+1}\}$ следва $\mathbb{E}_t\{Q_{t,t+1}e^{\,i_t}\} = 1$. Изваждайки двете, получаваме
\begin{equation}
\mathbb{E}_t \left\{ Q_{t,t+1} \left[ e^{\,i_t} - e^{\,i^i_t} \frac{\mathcal{E}_{i,t+1}}{\mathcal{E}_{i,t}} \right] \right\} = 0.
\label{eq:uip_exact}
\end{equation}
Лог-линеаризирането на \eqref{eq:uip_exact} около стационарно състояние при пълно предвиждане (в което $i_t = i^i_t$ и $\mathcal{E}_{i,t+1}/\mathcal{E}_{i,t} = 1$) и агрегирането по $i$ дават условието за непокрит лихвен паритет (uncovered interest parity)
\begin{equation}
i_t = i^*_t + \mathbb{E}_t\{ \Delta e_{t+1} \}.
\label{eq:uip}
\end{equation}

\paragraph{Условия на търговия и лихвени диференциали.}
Вземайки първа разлика на \eqref{eq:tot_exchange} и очаквания към момент $t$:
\[
\mathbb{E}_t\{ s_{t+1} \} - s_t
= \mathbb{E}_t\{ \Delta e_{t+1} \} + \mathbb{E}_t\{ \pi^*_{t+1} \} - \mathbb{E}_t\{ \pi_{H,t+1} \}.
\]
Заместването на $\mathbb{E}_t\{\Delta e_{t+1}\} = i_t - i^*_t$ от \eqref{eq:uip} и прегрупирането дават
\begin{equation}
s_t = \left( i^*_t - \mathbb{E}_t\{\pi^*_{t+1}\} \right) - \left( i_t - \mathbb{E}_t\{\pi_{H,t+1}\} \right) + \mathbb{E}_t\{ s_{t+1} \}.
\label{eq:tot_difference}
\end{equation}
Тъй като условията на търговия са еднозначно определени в стационарното състояние при пълно предвиждане, е в сила $\lim_{T \to \infty} \mathbb{E}_t\{ s_T \} = 0$ и \eqref{eq:tot_difference} може да бъде решено напред чрез рекурсивно заместване:
\begin{equation}
s_t = \mathbb{E}_t \left\{ \sum_{k=0}^{\infty} \left[ (i^*_{t+k} - \pi^*_{t+k+1}) - (i_{t+k} - \pi_{H,t+k+1}) \right] \right\}.
\label{eq:tot_forward}
\end{equation}
Уравненията \eqref{eq:tot_difference} и \eqref{eq:tot_forward} не представляват допълнително независимо равновесно условие: те се получават чрез комбиниране на уравненията на Ойлер за домашната и световната икономика с условието за споделяне на риска \eqref{eq:risk_sharing} и с връзката \eqref{eq:pi_cpi_domestic} \parencite[с.~158]{Gali2008}.

\section{Фирми}

В домашната икономика има континуум от фирми, индексирани с $j \in [0,1]$, функциониращи в условия на монополистична конкуренция. Всяка фирма произвежда точно една разновидност в количество $Y_t(j)$, при търсене, зададено с \eqref{eq:demand_variety}, и приема агрегатните величини за дадени.

\subsection{Технология и пределни разходи}

Типична фирма произвежда диференцирана стока с линейна технология
\begin{equation}
Y_t(j) = A_t N_t(j),
\label{eq:production}
\end{equation}
където $a_t := \ln A_t$ следва процес от вида AR(1):
\begin{equation}
a_t = \rho_a a_{t-1} + \varepsilon^a_t, \qquad \rho_a \in [0,1),
\label{eq:tech_process}
\end{equation}
а $\varepsilon^a_t$ е бял шум със средна нула.

Общите номинални разходи на фирмата са $TC^n_t(j) = (1-\tau) W_t N_t(j) = (1-\tau) W_t Y_t(j)/A_t$, където $\tau$ е субсидия за заетост, въведена с цел при подходящ избор на стойността ѝ разпределението при гъвкави цени да бъде ефективно. Номиналните пределни разходи са
\[
\psi_t = \frac{\partial TC^n_t(j)}{\partial Y_t(j)} = \frac{(1-\tau) W_t}{A_t},
\]
т.е. те не зависят от произвежданото количество и следователно са \emph{еднакви за всички домашни фирми}. Това е съществено опростяване спрямо базисния модел с намаляваща възвращаемост от мащаба. Реалните пределни разходи, изразени спрямо цените на домашните стоки, са $MC_t = \psi_t / P_{H,t}$, откъдето след логаритмуване
\begin{equation}
mc_t = -\nu + w_t - p_{H,t} - a_t,
\qquad \nu := -\ln(1-\tau).
\label{eq:mc_def}
\end{equation}

\subsection{Ценообразуване по схемата на Калво}

Приема се, че във всеки период дадена фирма с вероятност $1-\theta$ получава възможност да преоптимизира цената си, а с вероятност $\theta$ запазва цената от предходния период (схема на Калво, \cite{Calvo1983}). Вероятността за преоптимизиране не зависи от изминалото време от последната промяна. При допускането за симетрия всички фирми, които преоптимизират, решават една и съща задача и избират една и съща цена, означавана с $\bar{P}_{H,t}$.\footnote{В затворената икономика тази цена се означава с $P^*_t$. Тук звездичката е запазена за променливи, отнасящи се до световната икономика, поради което се използва $\bar{P}_{H,t}$.}

\subsubsection*{Необходимо условие за оптималност}

Фирмата, преоптимизираща в период $t$, решава задачата за максимизиране на условната печалба
\begin{equation}
\max_{\bar{P}_{H,t}} \ \mathbb{E}_t \sum_{k=0}^{\infty} \theta^k
\left\{ Q_{t,t+k} \left( \bar{P}_{H,t} Y_{t+k|t} - \Psi_{t+k}(Y_{t+k|t}) \right) \right\}
\label{eq:calvo_problem}
\end{equation}
при ограниченията
\begin{equation}
Y_{t+k|t} = \left( \frac{\bar{P}_{H,t}}{P_{H,t+k}} \right)^{-\varepsilon} Y_{t+k},
\qquad k = 0, 1, 2, \dots,
\label{eq:calvo_constraint}
\end{equation}
където $Y_{t+k}$ е агрегатното търсене на домашна продукция, а $\Psi_{t+k}(\cdot)$ е функцията на общите разходи.

Заместването на \eqref{eq:calvo_constraint} в \eqref{eq:calvo_problem} дава
\[
\mathbb{E}_t \sum_{k=0}^{\infty} \theta^k \left\{ Q_{t,t+k}
\left[ \bar{P}_{H,t} \left( \frac{\bar{P}_{H,t}}{P_{H,t+k}} \right)^{-\varepsilon} Y_{t+k}
- \Psi_{t+k}\!\left( \left( \frac{\bar{P}_{H,t}}{P_{H,t+k}} \right)^{-\varepsilon} Y_{t+k} \right) \right] \right\}.
\]
За получаване на необходимо условие за оптималност диференцираме по $\bar{P}_{H,t}$ и приравняваме на нула. Диференцирането на израза в квадратните скоби, при означение $\psi_{t+k|t} := \Psi'_{t+k}(Y_{t+k|t})$ за номиналните пределни разходи в период $t+k$ при цена, определена в период $t$, дава три члена:
\begin{equation}
\underbrace{\left( \frac{\bar{P}_{H,t}}{P_{H,t+k}} \right)^{-\varepsilon} Y_{t+k}}_{= \, Y_{t+k|t}}
+ \bar{P}_{H,t}(-\varepsilon) \bar{P}_{H,t}^{-\varepsilon-1} P_{H,t+k}^{\varepsilon} Y_{t+k}
- \psi_{t+k|t}(-\varepsilon) \bar{P}_{H,t}^{-\varepsilon-1} P_{H,t+k}^{\varepsilon} Y_{t+k}.
\label{eq:calvo_deriv}
\end{equation}
Тъй като $\bar{P}_{H,t}^{-\varepsilon-1} P_{H,t+k}^{\varepsilon} Y_{t+k} = Y_{t+k|t} / \bar{P}_{H,t}$, изразът \eqref{eq:calvo_deriv} се свежда до
\begin{equation}
Y_{t+k|t} - \varepsilon Y_{t+k|t} + \frac{\varepsilon}{\bar{P}_{H,t}} \psi_{t+k|t} Y_{t+k|t}
= Y_{t+k|t} \left( 1 - \varepsilon + \frac{\varepsilon}{\bar{P}_{H,t}} \psi_{t+k|t} \right).
\label{eq:calvo_collect}
\end{equation}
Изнасяйки $\frac{1-\varepsilon}{\bar{P}_{H,t}}$ пред скобата:
\begin{equation}
Y_{t+k|t} \left( 1 - \varepsilon + \frac{\varepsilon}{\bar{P}_{H,t}} \psi_{t+k|t} \right)
= \frac{(1-\varepsilon) Y_{t+k|t}}{\bar{P}_{H,t}}
\left( \bar{P}_{H,t} - \underbrace{\frac{\varepsilon}{\varepsilon-1}}_{=: \, \mathcal{M}} \psi_{t+k|t} \right).
\label{eq:calvo_factored}
\end{equation}

Приравнявайки очакваната дисконтирана сума на нула, използвайки линейното свойство на очакването и съкращавайки общия множител $\frac{1-\varepsilon}{\bar{P}_{H,t}} \ne 0$, получаваме необходимото условие за оптималност
\begin{equation}
\sum_{k=0}^{\infty} \theta^k\, \mathbb{E}_t \left\{ Q_{t,t+k} Y_{t+k|t}
\left[ \bar{P}_{H,t} - \mathcal{M} \psi_{t+k|t} \right] \right\} = 0.
\label{eq:foc_pricing}
\end{equation}
При гъвкави цени ($\theta = 0$) остава само членът с $k = 0$ и \eqref{eq:foc_pricing} се свежда до $\bar{P}_{H,t} = \mathcal{M}\psi_{t|t}$, което се интерпретира като формиране на цената чрез прилагане на надбавка $\mathcal{M}$ върху пределните разходи.

\subsubsection*{Лог-линеаризация на \eqref{eq:foc_pricing}}

Дефинираме ценовия индекс между периоди $t$ и $t+k$ като $\Pi_{H,t,t+k} := P_{H,t+k}/P_{H,t}$ и реалните пределни разходи в период $t+k$ при цена, определена в период $t$, като $MC_{t+k|t} := \psi_{t+k|t}/P_{H,t+k}$. Разделяйки \eqref{eq:foc_pricing} на $P_{H,t-1}$ и използвайки
\[
\frac{\psi_{t+k|t}}{P_{H,t-1}} = \frac{\psi_{t+k|t}}{P_{H,t+k}} \cdot \frac{P_{H,t+k}}{P_{H,t-1}} = MC_{t+k|t}\, \Pi_{H,t-1,t+k},
\]
получаваме
\begin{equation}
\sum_{k=0}^{\infty} \theta^k\, \mathbb{E}_t \left\{ Q_{t,t+k} Y_{t+k|t}
\left[ \frac{\bar{P}_{H,t}}{P_{H,t-1}} - \mathcal{M}\, MC_{t+k|t}\, \Pi_{H,t-1,t+k} \right] \right\} = 0.
\label{eq:foc_normalized}
\end{equation}

\paragraph{Стационарно състояние.}
В стационарно състояние с нулева инфлация имаме $\bar{P}_{H,t} = P_{H,t+k}$, следователно $\bar{P}_{H,t}/P_{H,t-1} = 1$ и $\Pi_{H,t-1,t+k} = 1$, откъдето $Y_{t+k|t} = Y$ и $MC_{t+k|t} = MC$. Освен това $Q_{t,t+k} = \beta^k$. Тогава \eqref{eq:foc_normalized} става
\[
Y \left[ 1 - \mathcal{M}\, MC \right] \sum_{k=0}^{\infty} (\beta\theta)^k = 0
\quad \Longrightarrow \quad
MC = \frac{1}{\mathcal{M}}
\quad \Longleftrightarrow \quad
mc = -\mu, \ \ \mu := \ln \mathcal{M}.
\]

\paragraph{Линеаризация.}
Записваме $Q_{t,t+k} = \beta^k \tilde{Q}_{t,t+k}$, където $\tilde{Q}_{t,t+k} := (C_{t+k}/C_t)^{-\sigma}(P_t/P_{t+k})$, и представяме всеки множител в експоненциален вид. Тогава \eqref{eq:foc_normalized} се разделя на две суми:
\begin{equation}
\sum_{k=0}^{\infty} (\beta\theta)^k \mathbb{E}_t \left\{ e^{\tilde{q}_{t,t+k}} e^{y_{t+k|t}} e^{\bar{p}_{H,t} - p_{H,t-1}} \right\}
= \sum_{k=0}^{\infty} (\beta\theta)^k \mathbb{E}_t \left\{ e^{\tilde{q}_{t,t+k}} e^{y_{t+k|t}} e^{\mu} e^{mc_{t+k|t}} e^{p_{H,t+k} - p_{H,t-1}} \right\}.
\label{eq:foc_exp}
\end{equation}
Означаваме стационарната стойност на израза под сумата отляво с $e^{EQ_1}$, а тази отдясно с $e^{EQ_2}$. От условието $MC = 1/\mathcal{M}$ следва $e^{EQ_1} = e^{EQ_2}$. Линеаризирайки двете страни около избраното стационарно състояние, получаваме
\begin{align}
&\sum_{k} (\beta\theta)^k \mathbb{E}_t \Big\{ e^{EQ_1} \big[ 1 + (\tilde{q}_{t,t+k} - \tilde{q}) + (y_{t+k|t} - y)
+ (\bar{p}_{H,t} - \bar{p}_H) - (p_{H,t-1} - p_H) \big] \Big\} \notag \\
= \ &\sum_{k} (\beta\theta)^k \mathbb{E}_t \Big\{ e^{EQ_2} \big[ 1 + (\tilde{q}_{t,t+k} - \tilde{q}) + (y_{t+k|t} - y)
+ \underbrace{(mc_{t+k|t} - mc)}_{=: \, \widehat{mc}_{t+k|t}} + (p_{H,t+k} - p_H) - (p_{H,t-1} - p_H) \big] \Big\}.
\label{eq:foc_linearized}
\end{align}
След съкращаване на $e^{EQ_1} = e^{EQ_2}$, на единиците, на членовете с $\tilde{q}_{t,t+k}$, на тези с $y_{t+k|t}$ и на тези с $p_{H,t-1}$ остава
\[
\sum_{k=0}^{\infty} (\beta\theta)^k \left( \bar{p}_{H,t} - p_{H,t-1} \right)
= \sum_{k=0}^{\infty} (\beta\theta)^k \mathbb{E}_t \left\{ \widehat{mc}_{t+k|t} + (p_{H,t+k} - p_{H,t-1}) \right\}.
\]
Лявата страна не зависи от $k$, поради което сумата е $(\bar{p}_{H,t} - p_{H,t-1})/(1-\beta\theta)$. Умножавайки по $(1-\beta\theta)$ и използвайки, че при линейна технология $\widehat{mc}_{t+k|t} = \widehat{mc}_{t+k}$ (пределните разходи не зависят от индивидуалната цена, съответно от индивидуалния обем на производството), получаваме
\begin{equation}
\bar{p}_{H,t} - p_{H,t-1} = (1-\beta\theta) \sum_{k=0}^{\infty} (\beta\theta)^k
\mathbb{E}_t \left\{ \widehat{mc}_{t+k} + (p_{H,t+k} - p_{H,t-1}) \right\}.
\label{eq:optimal_price_raw}
\end{equation}
Използвайки $\widehat{mc}_{t+k} = mc_{t+k} + \mu$ и $(1-\beta\theta)\sum_k(\beta\theta)^k = 1$, изразът \eqref{eq:optimal_price_raw} може да бъде записан и във вида
\begin{equation}
\bar{p}_{H,t} = \mu + (1-\beta\theta) \sum_{k=0}^{\infty} (\beta\theta)^k\, \mathbb{E}_t \left\{ mc_{t+k} + p_{H,t+k} \right\},
\label{eq:optimal_price_log}
\end{equation}
който се интерпретира като формиране на цената чрез надбавка над претеглените очаквани номинални пределни разходи \parencite[с.~159]{Gali2008}.

\subsubsection*{Рекурсивна форма}

За да се получи рекурсивен запис на \eqref{eq:optimal_price_raw}, разделяме дясната страна на две суми. Втората от тях се преобразува чрез $p_{H,t+k} - p_{H,t-1} = \sum_{m=0}^{k} \pi_{H,t+m}$ и смяна на реда на сумиране:
\[
(1-\beta\theta)\sum_{k=0}^{\infty}(\beta\theta)^k \sum_{m=0}^{k} \pi_{H,t+m}
= (1-\beta\theta)\sum_{m=0}^{\infty} \pi_{H,t+m} \sum_{k=m}^{\infty}(\beta\theta)^k
= \sum_{m=0}^{\infty}(\beta\theta)^m \pi_{H,t+m}.
\]
Следователно
\begin{equation}
\bar{p}_{H,t} - p_{H,t-1}
= \underbrace{(1-\beta\theta)\sum_{k}(\beta\theta)^k \mathbb{E}_t\{\widehat{mc}_{t+k}\}}_{=: \, \mathcal{A}_t}
+ \underbrace{\sum_{k}(\beta\theta)^k \mathbb{E}_t\{\pi_{H,t+k}\}}_{=: \, \mathcal{B}_t}.
\label{eq:price_two_sums}
\end{equation}
И двете суми имат рекурсивно представяне: $\mathcal{A}_t = (1-\beta\theta)\widehat{mc}_t + \beta\theta\,\mathbb{E}_t\{\mathcal{A}_{t+1}\}$ и $\mathcal{B}_t = \pi_{H,t} + \beta\theta\,\mathbb{E}_t\{\mathcal{B}_{t+1}\}$. Събирайки ги, получаваме
\begin{equation}
\bar{p}_{H,t} - p_{H,t-1} = \beta\theta\, \mathbb{E}_t\{ \bar{p}_{H,t+1} - p_{H,t} \}
+ (1-\beta\theta)\, \widehat{mc}_t + \pi_{H,t}.
\label{eq:price_recursive}
\end{equation}

\subsubsection*{Извеждане на общата динамика на цените}

Нека $\mathcal{J}(t) \subset [0,1]$ е множеството на фирмите, които не са получили възможност да актуализират цената си в период $t$. При симетрия всяка такава фирма има цена $P_{H,t-1}$, а останалите избират $\bar{P}_{H,t}$. Тогава общият ценови индекс има вида
\begin{equation}
P_{H,t} = \left[ \int_{\mathcal{J}(t)} P_{H,t-1}^{1-\varepsilon} dj + \int_{[0,1] \setminus \mathcal{J}(t)} \bar{P}_{H,t}^{1-\varepsilon} dj \right]^{\frac{1}{1-\varepsilon}}
= \left[ \theta P_{H,t-1}^{1-\varepsilon} + (1-\theta) \bar{P}_{H,t}^{1-\varepsilon} \right]^{\frac{1}{1-\varepsilon}},
\label{eq:price_index_dynamics}
\end{equation}
където второто равенство следва от (вариант на) закона за големите числа. При стандартната дефиниция $\Pi_{H,t} := P_{H,t}/P_{H,t-1}$, повдигайки \eqref{eq:price_index_dynamics} на степен $1-\varepsilon$ и разделяйки на $P_{H,t-1}^{1-\varepsilon}$, получаваме
\begin{equation}
\Pi_{H,t}^{1-\varepsilon} = \theta + (1-\theta) \left( \frac{\bar{P}_{H,t}}{P_{H,t-1}} \right)^{1-\varepsilon}.
\label{eq:pi_index}
\end{equation}
Лог-линеаризирайки \eqref{eq:pi_index} около стационарно състояние с нулева инфлация, т.е. $\Pi_{H,t} = \Pi_H = 1$ и $\bar{P}_{H,t} = \bar{P}_H = P_{H,t} = P_H$, лявата страна дава $(1-\varepsilon)\pi_{H,t}$, а дясната - $(1-\theta)(1-\varepsilon)(\bar{p}_{H,t} - p_{H,t-1})$. Съкращавайки $(1-\varepsilon)$:
\begin{equation}
\pi_{H,t} = (1-\theta)\left( \bar{p}_{H,t} - p_{H,t-1} \right).
\label{eq:pi_pstar}
\end{equation}

\subsubsection*{Уравнение за инфлацията}

Заместваме $\bar{p}_{H,t} - p_{H,t-1} = \frac{\pi_{H,t}}{1-\theta}$ от \eqref{eq:pi_pstar} в \eqref{eq:price_recursive}:
\[
\frac{\pi_{H,t}}{1-\theta} = \frac{\beta\theta}{1-\theta}\, \mathbb{E}_t\{\pi_{H,t+1}\} + (1-\beta\theta)\widehat{mc}_t + \pi_{H,t}.
\]
Прехвърляйки $\pi_{H,t}$ отляво и събирайки:
\[
\pi_{H,t}\left( \frac{1}{1-\theta} - 1 \right) = \pi_{H,t} \cdot \frac{\theta}{1-\theta}
= \frac{\beta\theta}{1-\theta}\, \mathbb{E}_t\{\pi_{H,t+1}\} + (1-\beta\theta)\widehat{mc}_t.
\]
Умножавайки по $\frac{1-\theta}{\theta}$, стигаме до уравнението за домашната инфлация
\begin{equation}
\pi_{H,t} = \beta\, \mathbb{E}_t\{ \pi_{H,t+1} \} + \lambda\, \widehat{mc}_t,
\qquad \lambda := \frac{(1-\beta\theta)(1-\theta)}{\theta}.
\label{eq:inflation_mc}
\end{equation}
Коефициентът $\lambda$ не зависи от нито един от параметрите, характеризиращи отворената икономика \parencite[с.~163]{Gali2008}. Той е и по-прост от съответния коефициент в базисния модел на затворената икономика, където при намаляваща възвращаемост в него влиза допълнителен множител, отчитащ зависимостта на пределните разходи от индивидуалния обем на производството. Тук този множител отпада, тъй като технологията \eqref{eq:production} е линейна.

\section{Равновесие}

\subsection{Пазар на стоки}

\subsubsection*{Търсене на домашна разновидност}

Условието за изравняване на търсенето и предлагането изисква производството на всяка домашна разновидност да покрива домашното търсене и търсенето от всички останали страни:
\begin{equation}
Y_t(j) = C_{H,t}(j) + \int_0^1 C^i_{H,t}(j)\, di.
\label{eq:market_clearing_j}
\end{equation}
Домашният компонент се получава от \eqref{eq:demand_variety} и \eqref{eq:demand_HF}:
\[
C_{H,t}(j) = \left( \frac{P_{H,t}(j)}{P_{H,t}} \right)^{-\varepsilon} (1-\alpha) \left( \frac{P_{H,t}}{P_t} \right)^{-\eta} C_t.
\]
При допускане за симетрични предпочитания в чужбина търсенето на страна $i$ за домашната разновидност $j$ преминава през същите три нива, но с ценови индекси, изразени във валутата на страна $i$:
\[
C^i_{H,t}(j) = \alpha \left( \frac{P_{H,t}(j)}{P_{H,t}} \right)^{-\varepsilon}
\left( \frac{P_{H,t}}{\mathcal{E}_{i,t} P^i_{F,t}} \right)^{-\gamma}
\left( \frac{P^i_{F,t}}{P^i_t} \right)^{-\eta} C^i_t.
\]
Заместването в \eqref{eq:market_clearing_j} дава
\begin{equation}
Y_t(j) = \left( \frac{P_{H,t}(j)}{P_{H,t}} \right)^{-\varepsilon}
\left[ (1-\alpha) \left( \frac{P_{H,t}}{P_t} \right)^{-\eta} C_t
+ \alpha \int_0^1 \left( \frac{P_{H,t}}{\mathcal{E}_{i,t} P^i_{F,t}} \right)^{-\gamma}
\left( \frac{P^i_{F,t}}{P^i_t} \right)^{-\eta} C^i_t\, di \right].
\label{eq:market_clearing}
\end{equation}

\subsubsection*{Агрегиране}

Агрегатното домашно производство се дефинира като
$Y_t := \left( \int_0^1 Y_t(j)^{1 - \frac{1}{\varepsilon}} dj \right)^{\frac{\varepsilon}{\varepsilon-1}}$.
Означавайки израза в квадратните скоби на \eqref{eq:market_clearing} с $\mathcal{D}_t$ (той не зависи от $j$), имаме
\[
\int_0^1 Y_t(j)^{\frac{\varepsilon-1}{\varepsilon}} dj
= \mathcal{D}_t^{\frac{\varepsilon-1}{\varepsilon}} P_{H,t}^{\varepsilon-1} \int_0^1 P_{H,t}(j)^{1-\varepsilon} dj
= \mathcal{D}_t^{\frac{\varepsilon-1}{\varepsilon}} P_{H,t}^{\varepsilon-1} P_{H,t}^{1-\varepsilon}
= \mathcal{D}_t^{\frac{\varepsilon-1}{\varepsilon}},
\]
където е използвана \eqref{eq:PH_index}. Повдигайки на степен $\frac{\varepsilon}{\varepsilon-1}$, получаваме $Y_t = \mathcal{D}_t$, т.е. агрегатното производство е точно изразът в квадратните скоби.

\subsubsection*{Преобразуване чрез относителни цени}

Изнасяме $(P_{H,t}/P_t)^{-\eta}$ пред скобата. За втория член вътре в скобата остава множителят
\[
\left( \frac{\mathcal{E}_{i,t} P^i_{F,t}}{P_{H,t}} \right)^{\gamma}
\left[ \frac{P^i_t}{P^i_{F,t}} \cdot \frac{P_{H,t}}{P_t} \right]^{\eta}
= \left( \frac{\mathcal{E}_{i,t} P^i_{F,t}}{P_{H,t}} \right)^{\gamma - \eta} \mathcal{Q}_{i,t}^{\eta},
\]
където второто равенство използва $\mathcal{Q}_{i,t} = \mathcal{E}_{i,t}P^i_t / P_t$. Освен това, от $P_{i,t} = \mathcal{E}_{i,t}P^i_{i,t}$, от $S_{i,t} = P_{i,t}/P_{H,t}$ и от дефиницията на ефективните условия на търговия за страна $i$, $S^i_t = P^i_{F,t}/P^i_{i,t}$, следва тъждеството
\begin{equation}
S^i_t\, S_{i,t}
= \frac{P^i_{F,t}}{P^i_{i,t}} \cdot \frac{\mathcal{E}_{i,t} P^i_{i,t}}{P_{H,t}}
= \frac{\mathcal{E}_{i,t} P^i_{F,t}}{P_{H,t}}.
\label{eq:tot_identity}
\end{equation}
Накрая, от условието за споделяне на риска \eqref{eq:risk_sharing_level} при $\vartheta_i = 1$ имаме $C^i_t = C_t\, \mathcal{Q}_{i,t}^{-1/\sigma}$, поради което $\mathcal{Q}_{i,t}^{\eta} C^i_t = C_t\, \mathcal{Q}_{i,t}^{\eta - 1/\sigma}$. Събирайки трите преобразувания:
\begin{equation}
Y_t = \left( \frac{P_{H,t}}{P_t} \right)^{-\eta} C_t
\left[ (1-\alpha) + \alpha \int_0^1 \left( S^i_t S_{i,t} \right)^{\gamma-\eta} \mathcal{Q}_{i,t}^{\eta - \frac{1}{\sigma}}\, di \right].
\label{eq:Y_exact}
\end{equation}

\subsubsection*{Лог-линеаризация на \eqref{eq:Y_exact}}

Логаритмувайки \eqref{eq:Y_exact}:
\[
y_t = -\eta\left( p_{H,t} - p_t \right) + c_t
+ \ln\!\left[ (1-\alpha) + \alpha \int_0^1 e^{(\gamma-\eta)(s^i_t + s_{i,t}) + (\eta - \frac{1}{\sigma}) q_{i,t}}\, di \right].
\]

\emph{Първи член.} От \eqref{eq:cpi_loglin} следва $p_{H,t} - p_t = -\alpha s_t$, поради което $-\eta(p_{H,t}-p_t) = \alpha\eta s_t$.

\emph{Трети член.} В симетричното стационарно състояние всички относителни цени са равни на единица, така че подинтегралната функция е $1$ и целият израз в скобите е $(1-\alpha) + \alpha = 1$. Приближението от първи ред на подинтегралната функция е
$1 + (\gamma-\eta)(s^i_t + s_{i,t}) + \left( \eta - \tfrac{1}{\sigma} \right) q_{i,t}$,
откъдето изразът в скобите става
\[
1 + \alpha \int_0^1 \left[ (\gamma - \eta)(s^i_t + s_{i,t}) + \left( \eta - \tfrac{1}{\sigma} \right) q_{i,t} \right] di.
\]
Използвайки $\int_0^1 s^i_t\, di = 0$ (сумата на условията на търговия по всички страни е нула), $\int_0^1 s_{i,t}\, di = s_t$ от \eqref{eq:tot_loglin} и $\int_0^1 q_{i,t}\, di = q_t$, а също и $\ln(1+x) \approx x$, третият член се свежда до
$\alpha(\gamma-\eta)s_t + \alpha\left( \eta - \tfrac{1}{\sigma} \right) q_t$.

\emph{Събиране.} Сумирайки трите приноса, членовете с $s_t$ се обединяват:
\begin{equation}
y_t = c_t + \underbrace{\alpha\eta s_t + \alpha(\gamma-\eta)s_t}_{= \, \alpha\gamma s_t}
+ \alpha\left( \eta - \frac{1}{\sigma} \right) q_t.
\label{eq:y_c_q}
\end{equation}
Заместването на $q_t = (1-\alpha)s_t$ от \eqref{eq:rer} и изнасянето на $\alpha/\sigma$ дават
\begin{equation}
y_t = c_t + \frac{\alpha}{\sigma}\Big[ \sigma\gamma + (1-\alpha)(\sigma\eta - 1) \Big] s_t
= c_t + \frac{\alpha\omega}{\sigma} s_t,
\qquad \omega := \sigma\gamma + (1-\alpha)(\sigma\eta - 1).
\label{eq:y_c_tot}
\end{equation}

\subsubsection*{Изчистване на световния пазар}

Аналогично условие е в сила за всяка страна: $y^i_t = c^i_t + \frac{\alpha\omega}{\sigma} s^i_t$. Интегрирайки по $i$ и използвайки отново $\int_0^1 s^i_t\, di = 0$, получаваме
\begin{equation}
y^*_t := \int_0^1 y^i_t\, di = \int_0^1 c^i_t\, di =: c^*_t.
\label{eq:world_clearing}
\end{equation}

Заместването на $c_t = c^*_t + \frac{1-\alpha}{\sigma}s_t$ от \eqref{eq:risk_sharing} в \eqref{eq:y_c_tot} и използването на \eqref{eq:world_clearing} дават
\[
y_t = y^*_t + \frac{1-\alpha}{\sigma}s_t + \frac{\alpha\omega}{\sigma}s_t
= y^*_t + \frac{(1-\alpha) + \alpha\omega}{\sigma}s_t
= y^*_t + \frac{1 + \alpha(\omega-1)}{\sigma}s_t,
\]
т.е.
\begin{equation}
y_t = y^*_t + \frac{1}{\sigma_\alpha} s_t,
\qquad \sigma_\alpha := \frac{\sigma}{1 + \alpha(\omega - 1)} > 0.
\label{eq:y_ystar_tot}
\end{equation}
Обръщайки \eqref{eq:y_ystar_tot}, получаваме връзката, която ще се използва многократно по-нататък:
\begin{equation}
s_t = \sigma_\alpha \left( y_t - y^*_t \right).
\label{eq:tot_output}
\end{equation}

\subsubsection*{Динамично уравнение за производството}

От \eqref{eq:y_c_tot} имаме $c_t = y_t - \frac{\alpha\omega}{\sigma}s_t$. Заместването в уравнението на Ойлер \eqref{eq:euler_log} дава
\begin{equation}
y_t - \frac{\alpha\omega}{\sigma}s_t
= \mathbb{E}_t\{y_{t+1}\} - \frac{\alpha\omega}{\sigma}\mathbb{E}_t\{s_{t+1}\}
- \frac{1}{\sigma}\left( i_t - \mathbb{E}_t\{\pi_{t+1}\} - \rho \right),
\label{eq:dis_step1}
\end{equation}
т.е.
\begin{equation}
y_t = \mathbb{E}_t\{ y_{t+1} \} - \frac{1}{\sigma} \left( i_t - \mathbb{E}_t\{\pi_{t+1}\} - \rho \right)
- \frac{\alpha\omega}{\sigma} \mathbb{E}_t\{ \Delta s_{t+1} \}.
\label{eq:dis_line1}
\end{equation}

\paragraph{Преминаване към домашна инфлация.}
От \eqref{eq:pi_cpi_domestic} следва $\mathbb{E}_t\{\pi_{t+1}\} = \mathbb{E}_t\{\pi_{H,t+1}\} + \alpha\,\mathbb{E}_t\{\Delta s_{t+1}\}$, поради което
\[
-\frac{1}{\sigma}\left( i_t - \mathbb{E}_t\{\pi_{t+1}\} - \rho \right)
= -\frac{1}{\sigma}\left( i_t - \mathbb{E}_t\{\pi_{H,t+1}\} - \rho \right)
+ \frac{\alpha}{\sigma}\mathbb{E}_t\{\Delta s_{t+1}\}.
\]
Заместването в \eqref{eq:dis_line1} и обединяването на членовете с $\mathbb{E}_t\{\Delta s_{t+1}\}$ дават
\begin{equation}
y_t = \mathbb{E}_t\{ y_{t+1} \} - \frac{1}{\sigma} \left( i_t - \mathbb{E}_t\{\pi_{H,t+1}\} - \rho \right)
- \frac{\alpha\Theta}{\sigma} \mathbb{E}_t\{ \Delta s_{t+1} \},
\label{eq:dis_line2}
\end{equation}
където\footnote{Означението $\Theta$ тук съответства на $\omega-1$ по \textcite[с.~161]{Gali2008} и не бива да се смесва с едноименния множител $\frac{1-\alpha}{1-\alpha+\alpha\varepsilon}$ от базисния модел на затворената икономика, където $\alpha$ обозначава степента на намаляваща възвращаемост. Тук последният отпада, тъй като технологията е линейна.}
\begin{equation}
\Theta := \omega - 1 = (\sigma\gamma - 1) + (1-\alpha)(\sigma\eta - 1).
\label{eq:Theta_def}
\end{equation}

\paragraph{Елиминиране на условията на търговия.}
От \eqref{eq:tot_output} следва $\mathbb{E}_t\{\Delta s_{t+1}\} = \sigma_\alpha\left( \mathbb{E}_t\{\Delta y_{t+1}\} - \mathbb{E}_t\{\Delta y^*_{t+1}\} \right)$. Означавайки $\chi := \frac{\alpha\Theta\sigma_\alpha}{\sigma}$ и замествайки в \eqref{eq:dis_line2}:
\[
y_t = \mathbb{E}_t\{y_{t+1}\} - \chi\left( \mathbb{E}_t\{y_{t+1}\} - y_t \right) + \chi\,\mathbb{E}_t\{\Delta y^*_{t+1}\}
- \frac{1}{\sigma}\left( i_t - \mathbb{E}_t\{\pi_{H,t+1}\} - \rho \right).
\]
Прехвърляйки членовете с $y_t$ отляво:
\[
(1-\chi)\, y_t = (1-\chi)\, \mathbb{E}_t\{y_{t+1}\} + \chi\,\mathbb{E}_t\{\Delta y^*_{t+1}\}
- \frac{1}{\sigma}\left( i_t - \mathbb{E}_t\{\pi_{H,t+1}\} - \rho \right).
\]
От дефиницията $\sigma_\alpha = \sigma/(1+\alpha\Theta)$ следва $\chi = \frac{\alpha\Theta}{1+\alpha\Theta}$ и
\[
1 - \chi = \frac{1}{1+\alpha\Theta} = \frac{\sigma_\alpha}{\sigma}.
\]
Разделяйки на $1-\chi = \sigma_\alpha/\sigma$ и използвайки $\frac{\sigma}{\sigma_\alpha}\chi = \frac{\sigma}{\sigma_\alpha}\cdot\frac{\alpha\Theta\sigma_\alpha}{\sigma} = \alpha\Theta$, получаваме окончателно
\begin{equation}
y_t = \mathbb{E}_t\{ y_{t+1} \} - \frac{1}{\sigma_\alpha} \left( i_t - \mathbb{E}_t\{\pi_{H,t+1}\} - \rho \right)
+ \alpha\Theta\, \mathbb{E}_t\{ \Delta y^*_{t+1} \}.
\label{eq:dis_levels}
\end{equation}
Степента на отвореност влияе върху чувствителността на производството към промени в домашния реален лихвен процент чрез $\sigma_\alpha$ \parencite[с.~161]{Gali2008}.

\subsection{Търговски баланс}

Нека $nx_t := \frac{1}{Y}\left( Y_t - \frac{P_t}{P_{H,t}} C_t \right)$ означава нетния износ в единици домашно производство, изразен като дял от стационарното производство $Y$. В симетричното стационарно състояние $Y = C$ и $P = P_H$, така че $nx = 0$. Приближението от първи ред е
\[
nx_t \approx \frac{1}{Y}\Big[ Y(1 + y_t) - C\big( 1 + (p_t - p_{H,t}) + c_t \big) \Big]
= y_t - c_t - \alpha s_t,
\]
където е използвано $Y = C$ и $p_t - p_{H,t} = \alpha s_t$. Заместването на $y_t - c_t = \frac{\alpha\omega}{\sigma}s_t$ от \eqref{eq:y_c_tot} дава
\begin{equation}
nx_t = \left( \frac{\alpha\omega}{\sigma} - \alpha \right) s_t
= \alpha \left( \frac{\omega}{\sigma} - 1 \right) s_t.
\label{eq:net_exports}
\end{equation}
Търговският баланс е балансиран във всеки момент, ако $\omega = \sigma$, т.е. ако $\sigma(\gamma-1) + (1-\alpha)(\sigma\eta-1) = 0$ \parencite[с.~162]{Gali2008}.

\subsection{Пазар на труд}

Изравняването на търсенето и предлагането на труд изисква $N_t := \int_0^1 N_t(j)\, dj$. Заместването на $N_t(j) = Y_t(j)/A_t$ от \eqref{eq:production} и на търсенето $Y_t(j) = (P_{H,t}(j)/P_{H,t})^{-\varepsilon}Y_t$ дава
\[
N_t = \frac{1}{A_t}\int_0^1 Y_t(j)\, dj
= \frac{Y_t}{A_t} \underbrace{\int_0^1 \left( \frac{P_{H,t}(j)}{P_{H,t}} \right)^{-\varepsilon} dj}_{=: \, d_t}.
\]
Величината $d_t$ може да се интерпретира като измерител на „разсейването“ (разпръснатостта) на цените между различните фирми. За нея може да се докаже, че в околност на стационарно състояние с нулева инфлация вариациите ѝ са от втори ред и следователно могат да бъдат изпуснати при приближение от първи ред. Логаритмувайки и полагайки $\ln d_t \approx 0$, получаваме линейното приближение за връзката между съвкупно производство, технологии и заетост
\begin{equation}
y_t = a_t + n_t.
\label{eq:agg_production}
\end{equation}

\subsection{Пределни разходи}

Изразът за реалните пределни разходи \eqref{eq:mc_def} се разгръща последователно:
\begin{align}
mc_t &= -\nu + (w_t - p_{H,t}) - a_t \notag \\
&= -\nu + \underbrace{(w_t - p_t)}_{\eqref{eq:labor_supply_log}} + \underbrace{(p_t - p_{H,t})}_{\eqref{eq:cpi_loglin}} - a_t \notag \\
&= -\nu + \sigma c_t + \varphi n_t + \alpha s_t - a_t.
\label{eq:mc_step}
\end{align}
От условието за споделяне на риска \eqref{eq:risk_sharing} и от \eqref{eq:world_clearing} следва
\[
\sigma c_t = \sigma c^*_t + (1-\alpha)s_t = \sigma y^*_t + (1-\alpha)s_t,
\]
а от агрегатната производствена функция \eqref{eq:agg_production} следва $n_t = y_t - a_t$. Заместването на двете в \eqref{eq:mc_step} дава
\begin{align}
mc_t &= -\nu + \sigma y^*_t + (1-\alpha)s_t + \varphi(y_t - a_t) + \alpha s_t - a_t \notag \\
&= -\nu + \sigma y^*_t + \varphi y_t + \underbrace{\big[(1-\alpha) + \alpha\big]}_{= \, 1} s_t - (1+\varphi) a_t \notag \\
&= -\nu + \sigma y^*_t + \varphi y_t + s_t - (1+\varphi) a_t.
\label{eq:mc_tot}
\end{align}
Заместването на $s_t = \sigma_\alpha(y_t - y^*_t)$ от \eqref{eq:tot_output} дава израз само чрез домашното производство, световното производство и технологията:
\begin{align}
mc_t &= -\nu + \sigma y^*_t + \varphi y_t + \sigma_\alpha y_t - \sigma_\alpha y^*_t - (1+\varphi)a_t \notag \\
&= -\nu + (\sigma_\alpha + \varphi) y_t + (\sigma - \sigma_\alpha) y^*_t - (1+\varphi) a_t.
\label{eq:mc_final}
\end{align}

От дефиницията $\sigma_\alpha = \sigma/(1+\alpha\Theta)$ следва $\sigma = \sigma_\alpha(1+\alpha\Theta)$, откъдето се получава тъждеството
\begin{equation}
\sigma - \sigma_\alpha = \alpha\Theta\,\sigma_\alpha,
\label{eq:sigma_identity}
\end{equation}
което ще бъде използвано няколкократно по-нататък. В частност, при $\alpha = 0$ и/или при $\Theta = 0$ е в сила $\sigma = \sigma_\alpha$ и домашните реални пределни разходи стават нечувствителни към движенията на чуждестранното производство \parencite[с.~164]{Gali2008}.

\subsection{Естествено ниво на производство}

Естественото ниво на производство $y^n_t$ се дефинира като нивото на производство в условия на гъвкави цени. Тогава реалните пределни разходи са постоянни и равни на $-\mu$, така че от \eqref{eq:mc_final}:
\[
-\mu = -\nu + (\sigma_\alpha + \varphi) y^n_t + (\sigma - \sigma_\alpha) y^*_t - (1+\varphi) a_t.
\]
Решавайки спрямо $y^n_t$:
\[
(\sigma_\alpha + \varphi) y^n_t = (\nu - \mu) + (1+\varphi)a_t - (\sigma - \sigma_\alpha) y^*_t,
\]
и разделяйки на $\sigma_\alpha + \varphi$, получаваме
\begin{equation}
y^n_t = \Gamma_0 + \Gamma_a a_t + \Gamma_* y^*_t,
\label{eq:natural_output}
\end{equation}
където, използвайки тъждеството \eqref{eq:sigma_identity} за третия коефициент,
\begin{equation}
\Gamma_0 := \frac{\nu - \mu}{\sigma_\alpha + \varphi},
\qquad
\Gamma_a := \frac{1+\varphi}{\sigma_\alpha + \varphi} > 0,
\qquad
\Gamma_* := -\frac{\sigma - \sigma_\alpha}{\sigma_\alpha + \varphi} = -\frac{\alpha\Theta\sigma_\alpha}{\sigma_\alpha + \varphi}.
\label{eq:Gamma_coefs}
\end{equation}

\section{Канонично представяне на равновесната динамика}

\subsection{Неокейнсианска крива на Филипс}

Нека $\tilde{y}_t := y_t - y^n_t$ означава отклонението от потенциалното производство (output gap). Изваждайки уравнението, дефиниращо $y^n_t$, от \eqref{eq:mc_final}, членовете с $y^*_t$ и $a_t$ се съкращават:
\[
mc_t - (-\mu) = (\sigma_\alpha + \varphi)(y_t - y^n_t),
\]
т.е.
\begin{equation}
\widehat{mc}_t = (\sigma_\alpha + \varphi)\, \tilde{y}_t.
\label{eq:mc_gap}
\end{equation}
Тук се използва, че $y^*_t$ е инвариантно спрямо събитията в домашната икономика.

Заместването на \eqref{eq:mc_gap} в \eqref{eq:inflation_mc} дава неокейнсианската крива на Филипс (New Keynesian Phillips Curve) за отворената икономика
\begin{equation}
\boxed{\ \pi_{H,t} = \beta\, \mathbb{E}_t\{ \pi_{H,t+1} \} + \kappa_\alpha \tilde{y}_t\ },
\qquad \kappa_\alpha := \lambda (\sigma_\alpha + \varphi).
\label{eq:nkpc}
\end{equation}
Формата на \eqref{eq:nkpc} съвпада с тази за затворената икономика; степента на отвореност влияе върху динамиката на домашната инфлация единствено чрез наклона $\kappa_\alpha$ \parencite[с.~164]{Gali2008}.

\subsection{Динамично IS уравнение}

Изваждаме $y^n_t$ от двете страни на \eqref{eq:dis_levels} и добавяме и изваждаме $\mathbb{E}_t\{y^n_{t+1}\}$ отдясно:
\begin{align}
\underbrace{y_t - y^n_t}_{\tilde{y}_t}
&= \underbrace{\mathbb{E}_t\{y_{t+1}\} - \mathbb{E}_t\{y^n_{t+1}\}}_{\mathbb{E}_t\{\tilde{y}_{t+1}\}}
+ \underbrace{\mathbb{E}_t\{y^n_{t+1}\} - y^n_t}_{\mathbb{E}_t\{\Delta y^n_{t+1}\}} \notag \\
&\quad - \frac{1}{\sigma_\alpha}\left( i_t - \mathbb{E}_t\{\pi_{H,t+1}\} - \rho \right)
+ \alpha\Theta\, \mathbb{E}_t\{ \Delta y^*_{t+1} \}.
\label{eq:dis_gap_raw}
\end{align}
Изнасяйки $-\frac{1}{\sigma_\alpha}$ пред всички членове, които не зависят от $\tilde{y}$, получаваме
\[
\tilde{y}_t = \mathbb{E}_t\{\tilde{y}_{t+1}\}
- \frac{1}{\sigma_\alpha}\Big[ i_t - \mathbb{E}_t\{\pi_{H,t+1}\}
- \underbrace{\big( \rho + \sigma_\alpha \mathbb{E}_t\{\Delta y^n_{t+1}\} + \sigma_\alpha\alpha\Theta\,\mathbb{E}_t\{\Delta y^*_{t+1}\} \big)}_{=: \ r^n_t} \Big],
\]
т.е. т.нар. динамично IS уравнение (динамична IS крива, dynamic IS equation)
\begin{equation}
\boxed{\ \tilde{y}_t = \mathbb{E}_t\{ \tilde{y}_{t+1} \} - \frac{1}{\sigma_\alpha} \left( i_t - \mathbb{E}_t\{ \pi_{H,t+1} \} - r^n_t \right)\ }
\label{eq:dis}
\end{equation}
където $r^n_t$ е естественото ниво на лихвения процент:
\begin{equation}
r^n_t := \rho + \sigma_\alpha \mathbb{E}_t\{\Delta y^n_{t+1}\} + \sigma_\alpha \alpha\Theta\, \mathbb{E}_t\{\Delta y^*_{t+1}\}.
\label{eq:rn_raw}
\end{equation}

\paragraph{Явен вид на естественото ниво на лихвения процент.}
Вземайки първа разлика на \eqref{eq:natural_output}, имаме
$\mathbb{E}_t\{\Delta y^n_{t+1}\} = \Gamma_a \mathbb{E}_t\{\Delta a_{t+1}\} + \Gamma_* \mathbb{E}_t\{\Delta y^*_{t+1}\}$.
От процеса \eqref{eq:tech_process} следва $\mathbb{E}_t\{a_{t+1}\} = \rho_a a_t$, откъдето
$\mathbb{E}_t\{\Delta a_{t+1}\} = (\rho_a - 1)a_t = -(1-\rho_a)a_t$.
Заместването в \eqref{eq:rn_raw} дава
\[
r^n_t = \rho - \sigma_\alpha \Gamma_a (1-\rho_a) a_t
+ \big( \sigma_\alpha \Gamma_* + \sigma_\alpha \alpha\Theta \big) \mathbb{E}_t\{\Delta y^*_{t+1}\}.
\]
Коефициентът пред $\mathbb{E}_t\{\Delta y^*_{t+1}\}$ се опростява чрез \eqref{eq:Gamma_coefs}:
\[
\sigma_\alpha \Gamma_* + \sigma_\alpha \alpha\Theta
= \sigma_\alpha \alpha\Theta \left( 1 - \frac{\sigma_\alpha}{\sigma_\alpha + \varphi} \right)
= \sigma_\alpha \alpha\Theta \cdot \frac{\varphi}{\sigma_\alpha + \varphi}.
\]
Следователно
\begin{equation}
r^n_t = \rho - \sigma_\alpha \Gamma_a (1-\rho_a)\, a_t
+ \frac{\alpha\Theta\sigma_\alpha\varphi}{\sigma_\alpha + \varphi}\, \mathbb{E}_t\{\Delta y^*_{t+1}\}.
\label{eq:rn}
\end{equation}
За разлика от затворената икономика тук $r^n_t \ne \rho + \sigma_\alpha\mathbb{E}_t\{\Delta y^n_{t+1}\}$: появява се допълнителен член, произхождащ от третия ред на \eqref{eq:dis_levels}, и именно комбинирането на двата приноса поражда коефициента $\frac{\alpha\Theta\sigma_\alpha\varphi}{\sigma_\alpha+\varphi}$.

Двете уравнения \eqref{eq:nkpc} и \eqref{eq:dis} формират неполитическия блок на модела. Той е формално идентичен с този за затворената икономика; отвореността се проявява само чрез коефициентите $\kappa_\alpha$ и $\sigma_\alpha$ и чрез структурата на $r^n_t$ \parencite[с.~165]{Gali2008}.

\section{Въвеждане на парична политика}

За да бъде решен моделът, е необходимо да се въведе правило за определяне на номиналния лихвен процент. Разглежда се правило на Тейлър от вида
\begin{equation}
i_t = \rho + \phi_\pi \pi_{H,t} + \phi_y \tilde{y}_t + v_t,
\qquad \phi_\pi, \phi_y \ge 0,
\label{eq:taylor}
\end{equation}
където $v_t$ е екзогенен компонент, следващ процес от вида AR(1):
\begin{equation}
v_t = \rho_v v_{t-1} + \varepsilon^v_t, \qquad \rho_v \in [0,1),
\label{eq:v_process}
\end{equation}
а $\varepsilon^v_t$ е бял шум със средна нула.

\subsection{Матрична форма на системата}

Заместваме \eqref{eq:taylor} в \eqref{eq:dis} и означаваме $\hat{r}^n_t := r^n_t - \rho$:
\[
\tilde{y}_t = \mathbb{E}_t\{\tilde{y}_{t+1}\}
- \frac{1}{\sigma_\alpha}\left( \phi_\pi \pi_{H,t} + \phi_y \tilde{y}_t + v_t - \mathbb{E}_t\{\pi_{H,t+1}\} - \hat{r}^n_t \right).
\]
Умножавайки по $\sigma_\alpha$ и прехвърляйки члена с $\phi_y \tilde{y}_t$ отляво:
\begin{equation}
(\sigma_\alpha + \phi_y)\tilde{y}_t = \sigma_\alpha \mathbb{E}_t\{\tilde{y}_{t+1}\}
+ \mathbb{E}_t\{\pi_{H,t+1}\} - \phi_\pi \pi_{H,t} + u_t,
\qquad u_t := \hat{r}^n_t - v_t.
\label{eq:system_step1}
\end{equation}
Заместваме $\pi_{H,t}$ от \eqref{eq:nkpc}:
\[
(\sigma_\alpha + \phi_y)\tilde{y}_t = \sigma_\alpha \mathbb{E}_t\{\tilde{y}_{t+1}\}
+ \mathbb{E}_t\{\pi_{H,t+1}\} - \phi_\pi \big( \beta\mathbb{E}_t\{\pi_{H,t+1}\} + \kappa_\alpha \tilde{y}_t \big) + u_t,
\]
откъдето след събиране на членовете с $\tilde{y}_t$
\begin{equation}
(\sigma_\alpha + \phi_y + \kappa_\alpha \phi_\pi)\tilde{y}_t
= \sigma_\alpha \mathbb{E}_t\{\tilde{y}_{t+1}\} + (1 - \beta\phi_\pi)\mathbb{E}_t\{\pi_{H,t+1}\} + u_t.
\label{eq:system_row1_raw}
\end{equation}
Полагайки
\begin{equation}
\Omega_\alpha := \frac{1}{\sigma_\alpha + \phi_y + \kappa_\alpha \phi_\pi},
\label{eq:Omega_def}
\end{equation}
първият ред на системата е
\begin{equation}
\tilde{y}_t = \Omega_\alpha \Big[ \sigma_\alpha \mathbb{E}_t\{\tilde{y}_{t+1}\} + (1-\beta\phi_\pi)\mathbb{E}_t\{\pi_{H,t+1}\} + u_t \Big].
\label{eq:system_row1}
\end{equation}

Заместването на \eqref{eq:system_row1} обратно в \eqref{eq:nkpc} дава
\[
\pi_{H,t} = \beta\mathbb{E}_t\{\pi_{H,t+1}\}
+ \kappa_\alpha \Omega_\alpha \Big[ \sigma_\alpha \mathbb{E}_t\{\tilde{y}_{t+1}\} + (1-\beta\phi_\pi)\mathbb{E}_t\{\pi_{H,t+1}\} + u_t \Big].
\]
Коефициентът пред $\mathbb{E}_t\{\pi_{H,t+1}\}$ се опростява, като се използва \eqref{eq:Omega_def}:
\begin{align}
\beta + \kappa_\alpha \Omega_\alpha (1-\beta\phi_\pi)
&= \beta\big( 1 - \kappa_\alpha \Omega_\alpha \phi_\pi \big) + \kappa_\alpha \Omega_\alpha \notag \\
&= \beta \Omega_\alpha \big( \sigma_\alpha + \phi_y + \kappa_\alpha\phi_\pi - \kappa_\alpha \phi_\pi \big) + \kappa_\alpha \Omega_\alpha \notag \\
&= \Omega_\alpha \Big[ \kappa_\alpha + \beta(\sigma_\alpha + \phi_y) \Big].
\label{eq:coef_simplify}
\end{align}

Обединявайки \eqref{eq:system_row1} и резултата за $\pi_{H,t}$, получаваме системата
\begin{equation}
\begin{bmatrix} \tilde{y}_t \\ \pi_{H,t} \end{bmatrix}
= A_\alpha
\begin{bmatrix} \mathbb{E}_t\{\tilde{y}_{t+1}\} \\ \mathbb{E}_t\{\pi_{H,t+1}\} \end{bmatrix}
+ B_\alpha \left( \hat{r}^n_t - v_t \right),
\label{eq:system}
\end{equation}
където
\begin{equation}
A_\alpha := \Omega_\alpha
\begin{bmatrix}
\sigma_\alpha & 1 - \beta\phi_\pi \\
\sigma_\alpha \kappa_\alpha & \kappa_\alpha + \beta(\sigma_\alpha + \phi_y)
\end{bmatrix},
\qquad
B_\alpha := \Omega_\alpha \begin{bmatrix} 1 \\ \kappa_\alpha \end{bmatrix}.
\label{eq:system_matrices}
\end{equation}

\subsection{Алтернативни правила}

В \textcite[с.~173]{Gali2008} се разглеждат три опростени правила за паричната политика, които ще бъдат използвани и в настоящото изследване:

\begin{itemize}
\item \emph{Правило на Тейлър спрямо домашната инфлация} (domestic inflation-based Taylor rule, DITR):
$i_t = \rho + \phi_\pi \pi_{H,t}$.
\item \emph{Правило на Тейлър спрямо ИПЦ инфлацията} (CPI inflation-based Taylor rule, CITR):
$i_t = \rho + \phi_\pi \pi_t$.
\item \emph{Фиксиран обменен курс} (exchange rate peg, PEG): $e_t = 0$ за всяко $t$.
\end{itemize}

Правилото DITR се получава от \eqref{eq:taylor} при $\phi_y = 0$ и $v_t \equiv 0$, поради което системата \eqref{eq:system} остава директно приложима с $\Omega_\alpha = 1/(\sigma_\alpha + \kappa_\alpha\phi_\pi)$. Правилата CITR и PEG изискват допълнителни уравнения, свързващи $\pi_t$ и $e_t$ с ендогенните променливи, чрез \eqref{eq:pi_cpi_domestic} и \eqref{eq:tot_exchange}.

\section{Екзогенни процеси}

Световното производство е екзогенно спрямо малката отворена икономика. Приема се, че то следва процес от вида AR(1)
\begin{equation}
y^*_t = \rho_{y^*} y^*_{t-1} + \varepsilon^{*}_t, \qquad \rho_{y^*} \in [0,1),
\label{eq:world_process}
\end{equation}
където $\varepsilon^{*}_t$ е бял шум със средна нула. Заедно с \eqref{eq:tech_process} и \eqref{eq:v_process} това напълно определя движещите сили на модела. Конкретните стойности на параметрите на трите процеса се въвеждат в глава~\ref{ch:experiment}.

\section{Обобщение}

Изведеният модел се свежда до следната система, която ще бъде използвана в останалата част от изложението:
\begin{align}
\pi_{H,t} &= \beta\, \mathbb{E}_t\{ \pi_{H,t+1} \} + \kappa_\alpha \tilde{y}_t, \tag{НККФ} \\
\tilde{y}_t &= \mathbb{E}_t\{ \tilde{y}_{t+1} \} - \frac{1}{\sigma_\alpha} \left( i_t - \mathbb{E}_t\{ \pi_{H,t+1} \} - r^n_t \right), \tag{ДИС} \\
i_t &= \rho + \phi_\pi \pi_{H,t} + \phi_y \tilde{y}_t + v_t, \tag{ПТ}
\end{align}
допълнена с процесите \eqref{eq:tech_process}, \eqref{eq:v_process} и \eqref{eq:world_process} за движещите сили и с идентичностите \eqref{eq:pi_cpi_domestic}, \eqref{eq:tot_exchange}, \eqref{eq:tot_output} и \eqref{eq:rer}, чрез които се възстановяват ИПЦ инфлацията, обменният курс, условията на търговия и реалният обменен курс.